\documentclass[11pt]{article}

\usepackage[T1]{fontenc}
\usepackage{amsmath,amssymb,amsthm}
\usepackage[margin=1in]{geometry}
\usepackage[hidelinks]{hyperref}

\newtheorem{theorem}{Theorem}
\newtheorem{lemma}[theorem]{Lemma}
\newtheorem{conjecture}[theorem]{Conjecture}
\newtheorem{proposition}[theorem]{Proposition}
\theoremstyle{definition}
\newtheorem{definition}[theorem]{Definition}
\theoremstyle{remark}
\newtheorem{remark}[theorem]{Remark}

\newcommand{\R}{\mathbb{R}}
\newcommand{\K}{\mathbf{K}}
\newcommand{\cK}{\mathcal{K}}
\newcommand{\eps}{\varepsilon}
\newcommand{\Lweak}{L^{3,\infty}}

\title{Type-II Record Transport: Residence, Physical Quantiles,\\
and Frozen-Band Correlation}
\author{John Robert Lawson and Codex}
\date{Research round ending 24 July 2026}

\begin{document}
\maketitle

\begin{abstract}
This report records the Type-II work completed after checkpoints
\texttt{868ba06} and \texttt{6ea929a}, together with the frozen-band
correlation theorem completed in the present commit batch.  The setting is
one hypothetical smooth finite-energy Navier--Stokes trajectory approaching
a first singular time, restricted to the energy-efficient carrier branch
with a positive terminal trace defect.  The new results are conditional
analytic reductions, not a proof of any Clay alternative.  They establish:
(i) logarithmic residence at a canonical defect scale; (ii) a physical
subgrid-quantile persistence-or-clock law between selected first records;
and (iii) actual old-band correlation or fixed nonlinear/viscous replacement
work.  Exact power and Hilbert-rotation ledgers show why the resulting
charges are still compatible with a Zeno cascade.
\end{abstract}

\section{Epistemic scope}

Throughout, \(u\) is a smooth finite-energy solution of
\[
 \partial_t u+\mathbb{P}\nabla\!\cdot(u\otimes u)=\nu\Delta u,
 \qquad \nabla\!\cdot u=0,
 \qquad 0\le t<T^*,
\]
where \(T^*\) is a putative first singular time.  Put
\[
 m(t):=\|u(t)\|_{\Lweak}.
\]
This report treats a subsequence of first record times \(t_j\uparrow T^*\)
for which \(m_j:=m(t_j)\to\infty\).  The carrier normalisation is
\[
 v_j(y,s):=\frac1{a_j}
 u(x_j+R_jy,t_j+\tau_js),
 \quad
 b_j:=a_j^2R_j^3,
 \quad
 \tau_j:=\frac{R_j}{a_j},
 \quad
 \eps_j:=\frac{\nu}{a_jR_j}.
\]
The energy-efficient branch gives, after a subsequence,
\[
 0<c\le b_j\le C,
 \qquad b_j\to b>0,
 \qquad
 \sup_{j,s}\|v_j(s)\|_2^2<\infty.
\]
The terminal carrier trace has a positive local defect measure
\(\mathcal{T}_0\): for some \(0\le\chi\le1\),
\[
 \Delta_\chi:=\int\chi\,d\mathcal{T}_0>0.
\]
Every theorem below is conditional on this branch.  No theorem below derives
the branch from arbitrary Clay data, controls the divergent-normalised-energy
branch, or proves global regularity or finite-time breakdown.

\section{First-record normalisation}

\begin{lemma}[Uniform critical ceiling]
The selected times may be chosen so that
\[
 m(t)\le m_j \qquad (0\le t\le t_j).
\]
For the half-extremising amplitude layer,
\[
 \left(\frac38\right)^{1/3}m_j
 \le a_jR_j\le m_j.
\]
Consequently
\[
 1\le\|v_j(0)\|_{\Lweak}\le\left(\frac83\right)^{1/3},
 \qquad
 \sup_{-t_j/\tau_j\le s\le0}
 \|v_j(s)\|_{\Lweak}
 \le\left(\frac83\right)^{1/3}.
\]
\end{lemma}

\begin{proof}
Smoothness implies continuity in \(L^2\cap L^\infty\), hence in
\(\Lweak\).  More explicitly, for \(0<\theta<1\), distribution-function
inclusion gives
\[
 \|f\|_{\Lweak}
 \le \frac{\|g\|_{\Lweak}}{1-\theta}
 +\frac{\|f-g\|_{\Lweak}}{\theta}.
\]
Applying this in both directions proves continuity of \(m(t)\).  Take
\(t_j\) to be the first hitting time of an increasing level \(M_j\to\infty\).

At \(t_j\), choose \(a_j\) so that
\[
 a_j^3|\{|u(t_j)|>a_j\}|\ge\frac12m_j^3
\]
and set
\[
 A_j:=\{a_j<|u(t_j)|\le2a_j\},
 \qquad R_j:=|A_j|^{1/3}.
\]
The weak-\(L^3\) upper bound at \(2a_j\) gives
\[
 |A_j|\ge
 \left(\frac12-\frac18\right)\frac{m_j^3}{a_j^3}
 =\frac38\frac{m_j^3}{a_j^3},
\]
while the upper bound at \(a_j\) gives
\[
 |A_j|\le\frac{m_j^3}{a_j^3}.
\]
Taking cube roots proves the amplitude--radius bounds.  Finally,
\[
 \|v_j(s)\|_{\Lweak}
 =\frac{m(t_j+\tau_js)}{a_jR_j},
\]
and the first-record property proves the asserted ceiling.
\end{proof}

\section{Global subgrid energy and endpoint flux}

Let
\[
 \Gamma_\ell(y):=(2\pi\ell^2)^{-3/2}
 \exp\!\left(-\frac{|y|^2}{2\ell^2}\right),
 \qquad
 \overline v_\ell:=\Gamma_\ell*v.
\]
Define
\[
 \cK_\ell[v]
 :=\frac12\left(\|v\|_2^2-\|\overline v_\ell\|_2^2\right)
 =\frac12\int_{\R^3}
 \left(1-e^{-\ell^2|\xi|^2}\right)|\widehat v(\xi)|^2\,d\xi.
\]
It is continuous and increasing in \(\ell\), and
\[
 \cK_\ell[v]\le\frac{\ell^2}{2}\|\nabla v\|_2^2.
\]

\begin{lemma}[Endpoint Lorentz flux ceiling]
If \(v\in L^2(\R^3)\cap\Lweak(\R^3)\), then
\[
 \left|\int_{\R^3}\Pi_\ell[v]\,dy\right|
 \le
 \frac{C\|v\|_{\Lweak}^3}{\ell}
 \log\!\left(
 e+\frac{C\|v\|_2^2}
 {\ell\|v\|_{\Lweak}^2}
 \right).
\]
In fact, the same right-hand side bounds
\(\|\Pi_\ell[v]\|_{L^1}\).
\end{lemma}

\begin{proof}
The subgrid stress satisfies
\[
 \|\tau_\ell(v,v)\|_{L^{3/2,\infty}}
 \le C\|v\|_{\Lweak}^2.
\]
For \(f=\nabla\overline v_\ell\), convolution gives
\[
 \|f\|_{L^{3,\infty}}\le\frac{CL}{\ell},
 \qquad
 \|f\|_\infty\le\frac{CL}{\ell^2},
 \qquad
 \|f\|_2\le\frac{C\sqrt E}{\ell},
\]
where \(L=\|v\|_{\Lweak}\) and \(E=\|v\|_2^2\).
If a function obeys
\[
 \|f\|_{L^{3,\infty}}\le A,\quad
 \|f\|_\infty\le B,\quad
 \|f\|_2\le C_2,
\]
then its decreasing rearrangement satisfies
\[
 f^*(r)\le\min\{B,Ar^{-1/3},C_2r^{-1/2}\}.
\]
Splitting the \(L^{3,1}\) integral at
\[
 r_0=(A/B)^3,\qquad r_1=(C_2/A)^6
\]
gives
\[
 \|f\|_{L^{3,1}}
 \le CA\left[
 1+\log_+\!\left(\frac{C_2^6B^3}{A^9}\right)
 \right].
\]
Substitution yields
\[
 \|\nabla\overline v_\ell\|_{L^{3,1}}
 \le
 \frac{CL}{\ell}\log\!\left(e+\frac{CE}{L^2\ell}\right).
\]
Lorentz H\"older applied pointwise to
\(\Pi_\ell=-\tau_\ell:\nabla\overline v_\ell\) completes the proof.
\end{proof}

\section{Canonical residence}

The exact global subgrid balance for carrier viscosity \(\eps_j\) is
\[
 \frac{d}{ds}\cK_\ell[v_j]
 =
 \int_{\R^3}\Pi_\ell[v_j]\,dy
 -
 \eps_j\int_{\R^3}d_\ell[v_j]\,dy,
 \qquad d_\ell[v_j]\ge0.
\]
Trace recovery gives
\[
 \lim_{\ell\downarrow0}\lim_{j\to\infty}
 \int\chi k_\ell(v_j(0);y)\,dy
 =\frac{\Delta_\chi}{2}.
\]
Fix \(0<\delta<\Delta_\chi/2\) and let \(\ell_j\) be the first
crossing satisfying
\[
 \cK_{\ell_j}[v_j(0)]=\delta.
\]
Then \(\ell_j\to0\).

\begin{theorem}[Logarithmic residence]
There is
\[
 h_j\gtrsim\frac{\ell_j}{\log(e+C/\ell_j)}
\]
such that
\[
 \cK_{\ell_j}[v_j(s)]\ge\frac{\delta}{2}
 \qquad(-h_j\le s\le0).
\]
Consequently
\[
 \frac{\eps_j}{\ell_j\log(e+C/\ell_j)}\longrightarrow0.
\]
\end{theorem}

\begin{proof}
The first-record ceiling, uniform carrier energy, and the endpoint flux lemma
give
\[
 \left|\int\Pi_{\ell_j}[v_j(s)]\right|
 \le
 B_j:=
 \frac{C}{\ell_j}\log(e+C/\ell_j)
\]
throughout the carrier past.  For \(s\le0\), the balance and
\(d_{\ell_j}\ge0\) imply
\[
 \delta-\cK_{\ell_j}[v_j(s)]
 \le B_j(0-s).
\]
Take \(h_j=\delta/(2B_j)\).  The Fourier inequality for \(\cK\) then gives
\[
 \|\nabla v_j(s)\|_2^2
 \ge\frac{\delta}{\ell_j^2}
 \qquad(-h_j\le s\le0).
\]
Therefore
\[
 \eps_j\int_{-h_j}^0\|\nabla v_j\|_2^2\,ds
 \gtrsim
 \frac{\eps_j}{\ell_j\log(e+C/\ell_j)}.
\]
The left side is bounded by the normalised pullback of the physical terminal
dissipation, which tends to zero by absolute continuity.  This proves the
claim.
\end{proof}

\section{Physical quantiles across records}

Define the physical subgrid energy
\[
 \K(r,t)
 :=\frac12\left(
 \|u(t)\|_2^2-\|\Gamma_r*u(t)\|_2^2
 \right).
\]
Gaussian scaling gives
\[
 \K(R_j\ell,t_j+\tau_js)
 =b_j\cK_\ell[v_j(s)].
\]
Choose two fixed physical levels
\[
 0<\eta_{\mathrm L}<\eta_{\mathrm H}<\frac{b\Delta_\chi}{2},
 \qquad
 \gamma:=\eta_{\mathrm H}-\eta_{\mathrm L}.
\]
Let \(r_j^\alpha\) be the first crossing
\[
 \K(r_j^\alpha,t_j)=\eta_\alpha,
 \qquad \alpha\in\{\mathrm L,\mathrm H\}.
\]
Then
\[
 \ell_j^\alpha:=\frac{r_j^\alpha}{R_j}\to0.
\]

\begin{theorem}[Cross-record quantile law]
Put
\[
 q_j:=\frac{m_{j+1}}{m_j},
 \qquad \Delta t_j:=t_{j+1}-t_j.
\]
For all sufficiently large \(j\), either
\[
 r_{j+1}^{\mathrm L}\le r_j^{\mathrm H},
\]
or
\[
 \Delta t_j
 \gtrsim
 \frac{\tau_j\ell_j^{\mathrm H}}
 {q_j^3\log(e+C/(q_j^2\ell_j^{\mathrm H}))}.
\]
\end{theorem}

\begin{proof}
The physical balance reads
\[
 \K(r,t_2)-\K(r,t_1)
 =
 \int_{t_1}^{t_2}\Phi(r,t)\,dt
 -
 \nu\int_{t_1}^{t_2}\!\!\int d_r[u]\,dx\,dt.
\]
The first-record property gives \(m(t)\le m_{j+1}\) on
\([t_j,t_{j+1}]\), hence
\[
 |\Phi(r,t)|
 \le
 F_j(r):=
 \frac{Cm_{j+1}^3}{r}
 \log\!\left(e+\frac{CE_0}{m_{j+1}^2r}\right).
\]
Let
\[
 D_j:=\nu\int_{t_j}^{T^*}\|\nabla u\|_2^2\,dt\to0.
\]
At \(r=r_j^{\mathrm H}\),
\[
 \eta_{\mathrm H}-\K(r_j^{\mathrm H},t_{j+1})
 \le F_j(r_j^{\mathrm H})\Delta t_j+D_j.
\]
If \(r_{j+1}^{\mathrm L}>r_j^{\mathrm H}\), first crossing gives
\(\K(r_j^{\mathrm H},t_{j+1})<\eta_{\mathrm L}\).  For large \(j\),
\(D_j\le\gamma/4\), and therefore
\[
 \Delta t_j\ge\frac{3\gamma}{4F_j(r_j^{\mathrm H})}.
\]
Using \(m_{j+1}=q_jm_j\), \(a_jR_j\asymp m_j\), and
\[
 \frac{b_j}{\tau_j}=a_j^3R_j^2
\]
converts this to the displayed carrier clock.
\end{proof}

\section{Frozen-band correlation}

Let \(G_rf=\Gamma_r*f\), and define
\[
 S_j:=G_{r_j^{\mathrm L}}^2-G_{r_j^{\mathrm H}}^2,
 \qquad B_j:=S_j^{1/2}.
\]
The Fourier multiplier of \(S_j\) is nonnegative:
\[
 e^{-(r_j^{\mathrm L})^2|\xi|^2}
 -
 e^{-(r_j^{\mathrm H})^2|\xi|^2}\ge0.
\]
Moreover,
\[
 \|B_ju(t_j)\|_2^2
 =2\bigl[\K(r_j^{\mathrm H},t_j)-\K(r_j^{\mathrm L},t_j)\bigr]
 =2\gamma.
\]

\begin{theorem}[Correlation or fixed replacement work]
Set
\[
 \psi_j:=S_ju(t_j),
 \qquad
 \mathcal C_j(t):=
 \langle u(t),\psi_j\rangle
 =\langle B_ju(t),B_ju(t_j)\rangle.
\]
At \(t_{j+1}\), at least one of the following holds:
\[
 \mathcal C_j(t_{j+1})\ge\gamma,
\]
\[
 \left|
 \int_{t_j}^{t_{j+1}}\!\!\int
 u\otimes u:\nabla\psi_j\,dx\,dt
 \right|\ge\frac{\gamma}{2},
\]
or
\[
 \left|
 \nu\int_{t_j}^{t_{j+1}}\!\!\int
 \nabla u:\nabla\psi_j\,dx\,dt
 \right|\ge\frac{\gamma}{2}.
\]
The first branch retains at least \(\gamma/4\) kinetic energy in the old
Gaussian band.  The second and third branches respectively force
\[
 \Delta t_j
 \gtrsim
 \frac{\tau_j\ell_j^{\mathrm L}}
 {q_j^2\log(e+C/\ell_j^{\mathrm L})}
\]
and
\[
 \Delta t_j
 \gtrsim
 \frac{\tau_j(\ell_j^{\mathrm L})^2}
 {\eps_jd_j},
 \qquad
 d_j:=\frac{D_j}{b_j}\to0.
\]
\end{theorem}

\begin{proof}
Pairing Navier--Stokes with the fixed divergence-free field \(\psi_j\)
gives the exact identity
\[
 \mathcal C_j(t)-2\gamma
 =
 \int_{t_j}^{t}\!\!\int u\otimes u:\nabla\psi_j
 -
 \nu\int_{t_j}^{t}\!\!\int\nabla u:\nabla\psi_j.
\]
If \(\mathcal C_j(t_{j+1})<\gamma\), the two work terms have a sum of
absolute magnitude greater than \(\gamma\); one therefore has magnitude at
least \(\gamma/2\).  In the correlation branch,
\[
 \gamma
 \le\|B_ju(t_{j+1})\|_2\sqrt{2\gamma},
\]
which proves the retained-energy constant.

Because \(G_r^2=G_{\sqrt2r}\), the endpoint gradient estimate gives
\[
 \|\nabla\psi_j\|_{L^{3,1}}
 \le
 \frac{Cm_j}{r_j^{\mathrm L}}
 \log\!\left(e+\frac{CE_0}{m_j^2r_j^{\mathrm L}}\right).
\]
Lorentz H\"older and \(m(t)\le q_jm_j\) yield
\[
 \left|\int u\otimes u:\nabla\psi_j\right|
 \le
 \frac{Cq_j^2}{\tau_j\ell_j^{\mathrm L}}
 \log(e+C/\ell_j^{\mathrm L}),
\]
which proves the nonlinear clock.

Also,
\[
 \|\nabla\psi_j\|_2\le\frac{C\sqrt{E_0}}{r_j^{\mathrm L}}.
\]
Thus
\[
 \left|
 \nu\int_{t_j}^{t_{j+1}}\!\!\int\nabla u:\nabla\psi_j
 \right|
 \le
 \frac{C\sqrt{\nu E_0D_j\Delta t_j}}{r_j^{\mathrm L}}.
\]
The fixed viscous-work floor and the exact identity
\[
 \frac{(r_j^{\mathrm L})^2}{\nu}
 =
 \frac{\tau_j(\ell_j^{\mathrm L})^2}{\eps_j}
\]
give the viscous clock.
\end{proof}

\section{Exact survivor ledger}

In units \(\nu=1\), take
\[
 m_j=2^{2j},\quad
 a_j=2^{6j},\quad
 R_j=2^{-4j},\quad
 \tau_j=2^{-10j},\quad
 \eps_j=2^{-2j},\quad q_j=4,
\]
and
\[
 \ell_j^{\mathrm H}=2^{-j},
 \qquad
 \ell_j^{\mathrm L}=2^{-j-1}.
\]
Then
\[
 \frac{\eps_j}
 {\ell_j^{\mathrm H}\log(e+C/\ell_j^{\mathrm H})}
 \asymp\frac{2^{-j}}j,
\]
\[
 \frac{\tau_j\ell_j^{\mathrm L}}
 {q_j^2\log(e+C/\ell_j^{\mathrm L})}
 \asymp\frac{2^{-11j}}j.
\]
Choose residence charges \(c_j\asymp2^{-j}/j\) and terminal tails
\[
 d_j:=\sum_{k>j}c_k\asymp\frac{2^{-j}}j.
\]
The viscous correlation clock is
\[
 \frac{\tau_j(\ell_j^{\mathrm L})^2}{\eps_jd_j}
 \asymp j2^{-9j}.
\]
All three series converge.  Taking
\(\Delta t_j\asymp2^{-11j}/j\) also gives
\[
 \sum_jm_{j+1}^4\Delta t_j
 \lesssim\sum_j\frac{2^{-3j}}j<\infty
\]
and
\[
 \frac{T^*-t_j}{\tau_j}\asymp\frac{2^{-j}}j\to0.
\]

For the correlation bookkeeping, let \((e_j)\) be orthonormal and rotate
\(\sqrt{2\gamma}\,e_j\) to \(\sqrt{2\gamma}\,e_{j+1}\) through angle
\(\pi/2\) on the \(j\)-th interval.  The squared norm stays \(2\gamma\), the
old correlation drops from \(2\gamma\) to zero, and the integrated work
against the old vector has magnitude \(2\gamma\).  An orthogonal reservoir
with initial energy exceeding \(\sum_jc_j\) pays the summable dissipation
charges.

\begin{remark}
This is exact Hilbert-space rotation and exact scalar bookkeeping.  It is not
a velocity field, it does not have the Navier--Stokes quadratic interaction,
and it is not a breakdown construction.  It proves only that all currently
derived inequalities remain mutually compatible.
\end{remark}

\section{Slam-dunk candidates, subject to review}

The following appear robust after full derivation and separate same-system
adversarial recomputation.  They still require independent mathematical
review before being treated as externally confirmed.

\begin{enumerate}
\item First-record sampling supplies a uniform weak-\(L^3\) ceiling over the
entire carrier past without imposing a Type-I bound on the original
trajectory.
\item A positive terminal trace defect has canonical physical
subgrid-energy quantiles whose normalised lengths tend to zero.
\item The endpoint Lorentz logarithm gives a genuine residence theorem and
the necessary inertial-range law
\[
 \frac{\eps_j}{\ell_j\log(e+C/\ell_j)}\to0.
\]
\item Consecutive selected records satisfy an exact physical
persistence-or-clock alternative.
\item The positive Gaussian band between two quantiles gives an exact
two-time correlation identity and a fixed replacement-work trichotomy.
\end{enumerate}

\section{Conjectures prompted by the round}

\begin{conjecture}[Common frozen-band budget]
For frozen Gaussian bands extracted from one Type-II first-record trajectory,
there is an adjoint or square-function construction for which the total
nonlinear replacement work is bounded by one finite physical quantity.
If true with an \(\ell^1\) event index, the fixed \(\gamma/2\) work floors
would exclude infinitely many replacements.
\end{conjecture}

\begin{conjecture}[Long-lived correlated-band accumulation]
After suitable thinning or recentering, either infinitely many old-band
correlations coexist at one physical time with bounded spectral overlap, or
their destruction pays a nonsummable charge.  Bessel orthogonality would
then contradict finite energy in the first alternative.
\end{conjecture}

\begin{conjecture}[NSE rotation obstruction]
The conservative orthogonal-band rotation ledger cannot be realised by the
Navier--Stokes triadic interaction under the first-record weak-\(L^3\)
ceiling and terminal dissipation collapse.
\end{conjecture}

\section{Open obligations}

\begin{enumerate}
\item Sum the event-dependent frozen-band nonlinear works against one common
same-trajectory budget, or construct a genuinely NSE-compatible survivor.
\item Recenter the global subgrid quantiles and correlated bands on the
retained amplitude layer.
\item In the coherent weak-trace branch, propagate a nonzero ancient Euler
object and prove a complete rigidity theorem.
\item Control all divergent-normalised-energy cells, where terminal physical
absolute continuity does not remove normalised dissipation.
\item Cover the energy-vanishing cells and the remaining geometry/clock
alternatives of the \(2\times3\times3\) Type-II ledger.
\item Keep ROUTE-R3B frozen pending genuinely new input or external review;
even its conditional closure would not cover Type II.
\end{enumerate}

\section{Repository and review record}

The proof-owning notes are
\begin{itemize}
\item \texttt{type-ii-record-flux-residence.md};
\item \texttt{type-ii-cross-record-quantile.md}; and
\item \texttt{type-ii-cross-record-correlation.md}.
\end{itemize}
Canonical experiment and route metadata remain in
\texttt{dossier/records/}.  This report is a round snapshot, not a competing
proof source.  The active public repository is
\url{https://github.com/Bingham-Research-Center/brc-navier-stokes}.

\bigskip
\noindent\textbf{Clay status:} unsolved.  No alternative A--D is proved.

\end{document}
